\chapter*{Abstract}
\addcontentsline{toc}{chapter}{Abstract}
\thispagestyle{plain}

Remote desktop access has become ordinary infrastructure for technical support,
distributed teams and unattended server administration. The tools in common use,
however, force an uncomfortable choice. Products such as TeamViewer and AnyDesk
are convenient but proprietary, and every session is mediated by vendor
infrastructure the user cannot inspect. Protocol-level alternatives such as VNC
and Microsoft RDP are open or well-documented but assume direct network
reachability, which is rarely available between two machines that both sit
behind network address translation.

This project presents \textbf{UpDesk}, a remote desktop system implemented
predominantly in Rust that addresses both concerns at once. Media and control
data travel over a direct WebRTC peer connection established between the
controlling and controlled machines, so screen frames never pass through a
central server in the common case. A small signaling server participates only
during connection establishment: it authenticates each endpoint using an
Ed25519 challenge--response handshake, relays session-description and ICE
candidate messages, and then steps aside. Because the server never possesses the
DTLS-SRTP keys negotiated by the peers, it cannot decrypt the stream it helped
to arrange.

The second contribution of the work is an \emph{unattended native host}. A
conventional web-technology host is constrained by the browser security model:
screen capture through \texttt{getDisplayMedia} demands a user gesture and a
picker dialogue, and it cannot observe the Windows login screen or User Account
Control prompts, which are rendered on an isolated secure desktop. UpDesk
therefore adds a standalone Rust binary that captures the framebuffer directly,
encodes it to H.264, publishes it on a native WebRTC track, and injects keyboard
and pointer events back into the active desktop. Registered as a Windows
service, this binary is available from boot and before any user has logged in.

The system was evaluated on a two-machine local network and across the public
Internet. Direct peer connectivity was achieved in the tested NAT
configurations, end-to-end latency remained in the range appropriate for
interactive use at 1080p, and the service-mode host survived reboot and session
switching without operator intervention. The report documents the protocol
design, the cryptographic reasoning behind it, the implementation of each
subsystem, and the measurements obtained.

\vspace{0.6cm}
\noindent\textbf{Keywords:} Remote Desktop, WebRTC, Peer-to-Peer, NAT Traversal,
Ed25519, Rust, Screen Capture, Windows Service, Tauri.

\clearpage
